\documentclass[11pt,a4paper]{amsart}
\usepackage{amssymb}
\usepackage[margin=1in]{geometry}
\usepackage[hidelinks]{hyperref}
\usepackage{array,longtable}
\setlength{\emergencystretch}{2em}

\newtheorem*{theorem}{Theorem}
\newtheorem{lemma}{Lemma}

\newcommand{\rem}{\operatorname{rem}}
\newcommand{\ValidSet}{\operatorname{ValidSet}}
\newcommand{\sFrac}{\operatorname{sFrac}}

\title{Disproof of Termination in Erd\H{o}s Problem 282 (Graham's Variant)}
\author{Ilia Pauzner}
\date{}

\begin{document}
\maketitle

For $A \subseteq \mathbb{N}$ and $x \in \mathbb{Q}$, define the greedy unit-fraction remainder sequence $\rem(A, x, t) \in \mathbb{Q}$ ($t \in \mathbb{N}$) recursively by $\rem(A, x, t) := 0$ when $x \le 0$, and for $x > 0$ by letting $n := \inf\{k \in \mathbb{N} : k \in A \text{ and } 1 / x \le k\}$, setting $r := x - 1 / n$, and putting
\[
\rem(A, x, 0) := r, \qquad \rem(A, x, t + 1) := \rem(A \setminus \{n\}, r, t).
\]
Graham~\cite{Gr64b} proved that a positive rational $m / n$ is a finite sum of distinct unit fractions with denominators $\equiv a \pmod{d}$ ($d > 1$) if and only if $\gcd\bigl(n / \gcd(n, a, d),\, d / \gcd(a, d)\bigr) = 1$. Erd\H{o}s and Graham~\cite[p.~31]{ErGr80} (see~\cite{EP282, FC282}) asked whether the greedy algorithm restricted to $\{n \in \mathbb{N} : n \equiv a \pmod{d}\}$ always terminates whenever this condition holds.

\begin{theorem}
For $x = 1 / 119 \in (0, 1)$, $a = 1$, and $d = 3$, Graham's coprimality condition holds ($\gcd(119, 3) = 1$), yet $\rem(\{n \in \mathbb{N} : n \equiv 1 \pmod{3}\}, 1 / 119, \cdot)$ is not eventually zero.
\end{theorem}

\subsection*{Correspondence with the Lean formalization}
The proof below is the one formalized in \nolinkurl{FormalConjectures/ErdosProblems/282.lean}, step by step, with the same names ($m$, $m'$, $q$, $q'$, $c$, $K$, $u_2$, $q_3$, $v_3$, $u_4$, $k_5$, $t$, $c'$, $T$) and with every identity and inequality in the form used in Lean. All variables are natural numbers unless noted otherwise. The correspondences are:
\begingroup
\small
\begin{longtable}{@{}>{\raggedright\arraybackslash}p{0.44\textwidth}>{\raggedright\arraybackslash}p{0.52\textwidth}@{}}
\hline
This note & Lean \\
\hline
\endhead
$\rem(A, x, t)$ & \texttt{greedyUnitFractionRem A x t} \\
$\ValidSet(A, q)$ & \texttt{ValidSet A q} \\
$\sFrac(m, q) := m / (3 m q - 1) \in \mathbb{Q}$ & \texttt{sFrac m q} \\
Lemma~\ref{lem:step} & \texttt{greedy\_step} \\
Lemma~\ref{lem:six} & \texttt{pos\_at\_six\_mul} \\
$u_2, q_3, v_3, u_4, k_5, t, c'$ & \texttt{u2}, \texttt{q3}, \texttt{v3}, \texttt{u4}, \texttt{k5}, \texttt{t}, \texttt{c'} \\
Theorem & \texttt{erdos\_282.variants.graham} \\
\hline
\end{longtable}
\endgroup

For $A \subseteq \mathbb{N}$ and $q \in \mathbb{N}$, define the predicate $\ValidSet(A, q)$ to mean that $1 \le q$ and for every $k \in \mathbb{N}$:
\[
(k \in A \Rightarrow k \bmod 3 = 1) \qquad\text{and}\qquad (3 q < k \text{ and } k \bmod 3 = 1 \Rightarrow k \in A).
\]
For $m, q \in \mathbb{N}$, put $\sFrac(m, q) := m / (3 m q - 1) \in \mathbb{Q}$.

\begin{lemma}\label{lem:step}
Let $m, m', q, q' \in \mathbb{N}$ and $A \subseteq \mathbb{N}$ satisfy $1 \le m$, $1 \le m'$, $\ValidSet(A, q)$, $q < q'$, and
\begin{equation}\label{eq:rel}
(m + 1)(3 m' q' - 1) = m' (3 m q - 1)(3 q + 1) \qquad \text{in } \mathbb{Q}.
\end{equation}
Then:
\begin{enumerate}
\item[(i)] $\ValidSet(A \setminus \{3 q + 1\}, q')$;
\item[(ii)] $\rem(A, \sFrac(m, q), 0) = \sFrac(m', q')$;
\item[(iii)] $\rem(A, \sFrac(m, q), t + 1) = \rem(A \setminus \{3 q + 1\}, \sFrac(m', q'), t)$ for all $t \in \mathbb{N}$.
\end{enumerate}
\end{lemma}

\begin{proof}
Since $1 \le m$ and $1 \le q$, we have $3 m q - 1 > 0$, so $\sFrac(m, q) > 0$. Set $S := \{k \in \mathbb{N} : k \in A \text{ and } 1 / \sFrac(m, q) \le k\}$. Because $3 q < 3 q + 1$ and $(3 q + 1) \bmod 3 = 1$, $\ValidSet(A, q)$ gives $3 q + 1 \in A$, and $3 m q - 1 \le (3 q + 1) m$ gives $1 / \sFrac(m, q) \le 3 q + 1$, so $3 q + 1 \in S$. Conversely, let $k \in S$. Then $3 m q \le k m + 1$. If $k < 3 q + 1$, then $k \bmod 3 = 1$ (from $k \in A$) forces $k + 2 \le 3 q$, so $(k + 2) m \le 3 m q \le k m + 1$, contradicting $1 \le m$. Hence $3 q + 1 \le k$ for all $k \in S$, which proves $\inf S = 3 q + 1$.

Next, clearing denominators against $(3 m q - 1)(3 q + 1) \ne 0$ and $3 m' q' - 1 \ne 0$, the identity $\sFrac(m, q) - 1 / (3 q + 1) = \sFrac(m', q')$ reduces to
\[
\bigl(m (3 q + 1) - (3 m q - 1)\bigr)(3 m' q' - 1) = m' (3 m q - 1)(3 q + 1),
\]
which holds by~\eqref{eq:rel} since $m (3 q + 1) - (3 m q - 1) = m + 1$. Finally, $\ValidSet(A \setminus \{3 q + 1\}, q')$ holds because $1 \le q < q'$ implies $1 \le q'$, membership in $A \setminus \{3 q + 1\}$ implies membership in $A$ (hence $k \bmod 3 = 1$), and any $k > 3 q'$ with $k \bmod 3 = 1$ satisfies $k > 3 q + 1$ (so $k \in A$ and $k \ne 3 q + 1$). Unfolding $\rem$ at $\sFrac(m, q) > 0$ now yields~(ii) and~(iii).
\end{proof}

\begin{lemma}\label{lem:six}
For every $K \in \mathbb{N}$, $A \subseteq \mathbb{N}$, and $c \in \mathbb{N}$ satisfying $\ValidSet(A, 2 c)$, $c \bmod 7 = 6$, and $(c \bmod 5 = 0 \text{ or } c \bmod 25 = 18)$, we have
\[
0 < \rem(A, \sFrac(1, 2 c), 6 K).
\]
\end{lemma}

\begin{proof}
We induct on $K$, quantifying over $A$ and $c$. Note first that $1 \le 2 c$ implies $1 \le c$.

\emph{Base case ($K = 0$).} Applying Lemma~\ref{lem:step}~(ii) with $(m, m', q, q') = (1, 2, 2 c, 6 c^2)$ (since $2 c < 6 c^2$ for $1 \le c$ and $2(36 c^2 - 1) = 2(6 c - 1)(6 c + 1)$) gives $\rem(A, \sFrac(1, 2 c), 0) = 2 / (36 c^2 - 1) > 0$.

\emph{Inductive step ($K \to K + 1$).} Set $u_2 := c^2 (36 c^2 + 1)$ and $q_3 := u_2 (9 u_2 + 1)$. Regardless of the parity of $u_2$, $u_2 (9 u_2 + 1)$ is even, so $q_3 \bmod 2 = 0$. Modulo $25$:
\begin{itemize}
\item if $c \bmod 5 = 0$, writing $c = 5 r$ gives $25 \mid c^2 \mid u_2 \mid q_3$, so $q_3 \bmod 25 = 0$;
\item if $c \bmod 25 = 18$, evaluating $u_2$ and $q_3$ modulo $25$ gives $q_3 \bmod 25 = 10$.
\end{itemize}
In either case $10 \mid q_3$, so write $q_3 = 10 v_3$. Set $u_4 := v_3 (40 v_3 + 1)$ and $k_5 := u_4 (45 u_4 + 2)$. From $c \bmod 7 = 6$ we compute $q_3 \bmod 7 = 3$, hence $v_3 \bmod 7 = 1$ (since $10 v_3 = q_3$), hence $u_4 \bmod 7 = 6$ and $k_5 \bmod 7 = 1$. Writing $t := \lfloor k_5 / 7 \rfloor$ (so $k_5 = 7 t + 1$) and setting
\[
c' := 7 (36 t^2 + 11 t) + 6,
\]
ring expansion gives the exact identity
\begin{equation}\label{eq:7c}
7 c' = 36 k_5^2 + 5 k_5 + 1,
\end{equation}
and visibly $c' \bmod 7 = 6$. For the modulo-$5$/$25$ disjunction on $c'$, note from $10 v_3 = q_3$ and $q_3 \bmod 25 \in \{0, 10\}$ that $v_3 \bmod 5 \in \{0, 1\}$:
\begin{itemize}
\item if $v_3 \bmod 5 = 0$, then $5 \mid v_3 \mid u_4 \mid k_5$; writing $k_5 = 5 r$ in~\eqref{eq:7c} gives $7 c' = 25(36 r^2 + r) + 1$, whence $c' \bmod 25 = 18$;
\item if $v_3 \bmod 5 = 1$, then $u_4 \bmod 5 = 1$ and $k_5 \bmod 5 = 2$, so~\eqref{eq:7c} gives $(7 c') \bmod 5 = (36 \cdot 4 + 10 + 1) \bmod 5 = 0$, whence $c' \bmod 5 = 0$.
\end{itemize}

Now apply Lemma~\ref{lem:step} six times starting from $A_0 := A$:
\begin{enumerate}
\item $(m, m', q, q') = (1, 2, 2 c, 6 c^2)$: $2 c < 6 c^2$ since $1 \le c$, and $2(36 c^2 - 1) = 2(6 c - 1)(6 c + 1)$;
\item $(m, m', q, q') = (2, 3, 6 c^2, 2 u_2)$: $6 c^2 < 2 c^2 (36 c^2 + 1)$ since $1 \le c^2$, and $3(18 u_2 - 1) = 3(36 c^2 - 1)(18 c^2 + 1)$;
\item $(m, m', q, q') = (3, 4, 2 u_2, q_3)$: $2 u_2 < u_2 (9 u_2 + 1)$ since $1 \le 2 u_2$, and $4(12 q_3 - 1) = 4(18 u_2 - 1)(6 u_2 + 1)$;
\item $(m, m', q, q') = (4, 5, q_3, 6 u_4)$: using $q_3 = 10 v_3$, $10 v_3 < 6 v_3 (40 v_3 + 1)$ since $1 \le q_3$, and $5(90 u_4 - 1) = 5(120 v_3 - 1)(30 v_3 + 1)$;
\item $(m, m', q, q') = (5, 6, 6 u_4, 2 k_5)$: $6 u_4 < 2 u_4 (45 u_4 + 2)$ since $1 \le 6 u_4$, and $6(36 k_5 - 1) = 6(90 u_4 - 1)(18 u_4 + 1)$;
\item $(m, m', q, q') = (6, 1, 2 k_5, 2 c')$: using~\eqref{eq:7c}, $2 k_5 < 2 c'$ since $1 \le 2 k_5$, and $7(6 c' - 1) = (36 k_5 - 1)(6 k_5 + 1)$ since $7(6 c' - 1) = 6(7 c') - 7 = 6(36 k_5^2 + 5 k_5 + 1) - 7 = 216 k_5^2 + 30 k_5 - 1 = (36 k_5 - 1)(6 k_5 + 1)$.
\end{enumerate}
Writing $6(K + 1) = 6 K + 1 + 1 + 1 + 1 + 1 + 1$, the six step equations~(iii) rewrite the remainder $\rem(A, \sFrac(1, 2 c), 6(K + 1))$ to $\rem(A_6, \sFrac(1, 2 c'), 6 K)$ with $\ValidSet(A_6, 2 c')$, which is positive by the induction hypothesis applied to $A_6$ and $c'$.
\end{proof}

\begin{proof}[Proof of the Theorem]
Take $x = 1 / 119 \in (0, 1)$, $a = 1$, and $d = 3$; then $\gcd(1, 3) = 1$ and $\gcd(119, 3) = 1$. Suppose for contradiction that $\rem(A, 1 / 119, t) = 0$ for all $t \ge T$, where $A := \{n \in \mathbb{N} : n \equiv 1 \pmod{3}\}$. For $c = 20$, we have $\ValidSet(A, 40)$, $20 \bmod 7 = 6$, $20 \bmod 5 = 0$, and $\sFrac(1, 40) = 1 / 119$. Lemma~\ref{lem:six} with $K = T$ gives $0 < \rem(A, 1 / 119, 6 T)$, contradicting $\rem(A, 1 / 119, 6 T) = 0$ since $T \le 6 T$.
\end{proof}

\begin{thebibliography}{9}
\bibitem{Gr64b}
R.~L.~Graham, \emph{On finite sums of unit fractions}, Proc.\ London Math.\ Soc.\ (3) \textbf{14} (1964), 193--207.

\bibitem{ErGr80}
P.~Erd\H{o}s and R.~L.~Graham, \emph{Old and new problems and results in combinatorial number theory}, Monographies de L'Enseignement Math\'ematique, vol.~28, Universit\'e de Gen\`eve, 1980.

\bibitem{EP282}
\emph{Erd\H{o}s Problem \#282}, \url{https://www.erdosproblems.com/282}.

\bibitem{FC282}
The Formal Conjectures Authors, \emph{FormalConjectures/ErdosProblems/282.lean},\\ \url{https://github.com/google-deepmind/formal-conjectures}.
\end{thebibliography}

\end{document}
